\documentclass{article}
\usepackage{../../Self_Style}

\title{Phys 115B HW 9}
\author{Zih-Yu Hsieh}
\date{\today}

\begin{document}
\maketitle

\section*{1 D}
\begin{ques}\label{q1}
Consider the parity operator $\hat{\Pi}$ in three dimensions, such that $\hat{\Pi}\psi(\mathbf r)=\psi(-\mathbf r)$.

\begin{itemize}
\item[(a)] Show that for states $\psi$ expressed in polar coordinates, the action of the parity operator can be written as
\[
\hat{\Pi}\psi(r,\theta,\phi)=\psi(r,\pi-\theta,\phi+\pi)
\]

\item[(b)] Show that for hydrogenic energy eigenstates $\psi_{nlm}(r,\theta,\phi)$ we have
\[
\hat{\Pi}\psi_{nlm}=(-1)^l\psi_{nlm}(r,\theta,\phi)
\]
\end{itemize}
\end{ques}

\begin{proof}
    
    \hfil

    \begin{itemize}
        \item[(a)] Given the parity operator has the coordinate change $\br \mapsto -\br$, one has the following on each component:
        \begin{align}
            &x: \quad r\sin(\theta)\cos(phi) \mapsto -r\sin(\theta)\cos(\phi) = r\sin(\pi-\theta)\cos(\phi+\pi)\\
            &y: \quad r\sin(\theta)\sin(phi) \mapsto -r\sin(\theta)\sin(\phi) = -r\sin(\pi-\theta)\sin(\phi) = r\sin(\theta)\sin(\phi+\pi)\\
            &z:\quad r\cos(theta) \mapsto -r\cos(\theta) = r\cos(\pi-\theta)
        \end{align}
        So, one transformation on the spherical coordinates that's compatible with the parity operator is $r\mapsto r$, $\theta\mapsto \pi-\theta$, and $\phi\mapsto \phi+\pi$.

        Hence, using $r,\theta,\psi$ as the input, one has:
        \begin{align}
            \hat{\Pi}\psi(r,\theta,\phi) = \psi(r,\pi-\theta, \phi+\pi)
        \end{align}

        \hfil

        \item[(b)] Since for $\psi_{nlm} = R_{n,l} (r) Y^m_l (\theta,\phi)$ in general, then the only one affected by the parity is $Y^m_l(\theta,\phi)$. 
        
        Recall the formula $Y^m_l(\theta,\phi) = e^{im\phi}P^m_l (\cos(\theta))$, where $P^m_l$ is given as follow:
        \begin{align}
            P^m_l (z)=(1-z^2)^{|m|/2}\left(\frac{d}{dz}\right)^{|m|}P_l (z), \quad P_l (z) := \frac{1}{2^l l!}\left(\frac{d}{dz}\right)^l (z^2-1)^l
        \end{align}
    \end{itemize}
    Which, it can be shown inductively, that all $P_l (z)$ only contains the monomials with the degree of same parity with $l$; then, one can also inductively shown $\left(\frac{d}{dz}\right)^{|m|}P_l (z)$ has all monomial with degree of same parity with $l-|m|$, showing $P^m_l (-z) = (-1)^{l-|m|}P^m_l (z)$. Finally, one has $e^{im(\phi+\pi)} = (-1)^{|m|}e^{im\phi}$. Hence, with $\cos(\pi-\theta) = -\cos(\theta)$, one has the following:
    \begin{align}
        \hat{\Pi}\psi_{nlm} &= \psi_{nlm}(r,\pi-\theta, \phi+\pi) = R_{n,l}(r)e^{im(\phi+\pi)}P^m_l (\cos(\pi-\theta))\\
        &= R_{n,l}(r)(-1)^{|m|}e^{im\phi}P^m_l (-cos(\theta))= R_{n,l} (-1)^{|m|}e^{im\phi}(-1)^{l-|m|}P^m_l (\cos(\theta))\\
        &= (-1)^l R_{n,l}e^{im\phi}P^m_l (\cos(\theta))= (-1)^l\psi_{nlm}(r,\theta,\phi)
    \end{align} 
\end{proof}

\newpage

\section*{2}
\begin{ques}\label{q2}
We can classify operators depending on how they transform under parity, $\hat{Q}\rightarrow \hat{\Pi}^\dagger \hat{Q}\hat{\Pi}$ (see example 6.2 or Table 6.1 in the book can be useful).

\begin{itemize}
\item[(a)] In one dimension, show that position and momentum are odd under parity.

\item[(b)] In three dimensions, operators that are invariant under parity and commute with $\hat L$ are called scalar operators, $\hat{\Pi}^\dagger \hat Q_s \hat{\Pi}=\hat Q_s$. Operators that flip sign under parity and commute with $\hat L$ are called pseudo-scalar operators, $\hat{\Pi}^\dagger \hat Q_s \hat{\Pi}=-\hat Q_s$. Show that scalar operators satisfy $[\hat{\Pi},\hat Q_s]=0$, while pseudoscalar operators satisfy
\[
\{\hat{\Pi},\hat Q_s\}=0
\]
(where the curly brackets denote the anti-commutator).

\item[(c)] Operators that flip sign under parity and whose commutators with $\hat L$ are like those of $\hat L$ are called vector operators, such that $\hat{\Pi}^\dagger \hat Q_v \hat{\Pi}=-\hat Q_v$. Operators that are invariant under parity and whose commutators with $\hat L$ are like those of $\hat L$ are called pseudovector or axial vector operators, $\hat{\Pi}^\dagger \hat Q_a \hat{\Pi}=\hat Q_a$. Show that vector operators satisfy
\[
\{\hat{\Pi},\hat Q_v\}=0
\]
and pseudo-vectors satisfy
\[
[\hat{\Pi},\hat Q_a]=0.
\]

\item[(d)] Give an example of a scalar operator, pseudoscalar operator, a vector operator and an axial or pseudo vector operator (feel free to Google or use the book).
\end{itemize}
\end{ques}

\begin{proof}

    \hfil

    \begin{itemize}
        \item[(a)] Let's first recall some tools, for instance $\braket{x}{p}=e^{ipx/\hbar}$. Then, we have the following:
        \begin{align}
            \hat{Pi}^\dagger \hat{x} \hat{\Pi} \ket{x} = \hat{\Pi}\hat{x}\ket{-x} = -x\hat{Pi}\ket{-x} = -x\ket{x}
        \end{align}
        Since above is true for all $x\in\RR$, it shows $\hat{\Pi}^\dagger\hat{x}\hat{\Pi}=-\hat{x}$.

        Similarly, we have the following:
        \begin{align}
            \hat{\Pi}\ket{p} &= \int \hat{\Pi}\ket{x}\braket{x}{p}dx = \int e^{ipx/\hbar}\ket{-x}dx\\
            &= \int e^{i(-p)u/\hbar}\ket{u}du = \ket{-p}
        \end{align}
        (Note: Here $u:=-x$).

        Hence, the full relation on momentum is:
        \begin{align}
            \hat{\Pi}^\dagger\hat{p}\hat{\Pi}\ket{p} = \hat{\Pi}\hat{p}\ket{-p} = -p\hat{\Pi}\ket{-p} = -p\ket{p}
        \end{align}
        Since above is true for all $p \in\RR$, it shows $\hat{\Pi}^\dagger\hat{p}\hat{\Pi}=-\hat{p}$.

        \hfil

        \item[(b)] For scalar operators we have the following:
        \begin{align}
            \hat{\Pi}^\dagger \hat{Q}_s \hat{\Pi} = \hat{Q}_s \implies \hat{Q}_s \hat{\Pi} = \hat{\Pi}\hat{Q}_s \implies [\hat{\Pi},\hat{Q}_s]=0
        \end{align}
        For pseudo-scalar operators we have the following:
        \begin{align}
            \hat{\Pi}^\dagger\hat{Q}_s\hat{\Pi}=-\hat{Q}_s \implies \hat{Q}_s \hat{\Pi}=-\hat{\Pi}\hat{Q}_s \implies \{\hat{\Pi},\hat{Q}_s\}=0
        \end{align}

        \hfil

        \item[(c)]
    \end{itemize}
\end{proof}

\newpage

\section*{3}
\begin{ques}\label{q3}
Consider a particle of mass $m$ and charge $q$ confined to a circular ring in the $x$-$y$ plane, with radius $r$. The system will be driven by an electric field $E_x$ in the $x$ direction (which can stimulate photon induced transitions). It is natural to work in polar coordinates for this problem, with $x=r\cos\theta$.

\begin{itemize}
\item[(a)] What are the first five energies of the system with $E_x=0$?

\item[(b)] What are the first five eigenfunctions $\psi(\theta)$? Please normalize them, and pick your basis such that the eigenfunctions are also eigenfunctions of $L_z$.

\item[(c)] Explicitly write out a $3\times3$ matrix for the Hamiltonian
\[
H' = E_x q x = E q r \cos\theta
\]
This matrix will use the three lowest states above. Extrapolate this result to write a concise selection rule that tells what transitions are allowed for any state.

\item[(d)] What frequencies are absorbed by a “gas” of such rings if they are cold enough to be initially in the ground state?

\item[(e)] What additional frequencies are absorbed by rings initially in the first excited state?
\end{itemize}
\end{ques}

\begin{proof}
\end{proof}

\newpage

\end{document}